\begin{ouverture}[courbe]{4}{Développer, factoriser et prouver}{Partie 1 · Nombres et calculs}
{Expressions · Distributivité · Identités · Programmes de calcul}
\objectif{Transformer une expression sans changer sa valeur.}
\objectif{Choisir entre développement et factorisation.}
\objectif{Démontrer avec une lettre.}
\prerequis{Produits, sommes et calculs avec les relatifs.}
\end{ouverture}
\pagestyle{otMain}\essentiel\fontsize{11}{15}\selectfont
\essSection{1}{Lire et développer}
\begin{definition}Une expression littérale décrit un calcul avec des lettres. Réduire consiste à regrouper des termes semblables. Développer transforme un produit en somme.\end{definition}
\begin{propriete}La distributivité donne $k(a+b)=ka+kb$. En particulier, $-(a-b)=-a+b$. Pour deux parenthèses :
\[(a+b)(c+d)=ac+ad+bc+bd.\]\end{propriete}
\begin{exemple}$(2x-3)(x+4)=2x^2+8x-3x-12=2x^2+5x-12$. Le terme $2x^2$ est un carré multiplié par 2, pas un terme en $x$.\end{exemple}
\begin{savoirfaire}[Remplacer une lettre]Pour $x=-2$, écrire les parenthèses : $3x^2-5x=3(-2)^2-5(-2)=22$. Calculer les puissances avant les produits.\end{savoirfaire}
\begin{attention}$3x+2x=5x$, mais $3x+2$ ne se réduit pas à $5x$. Un test numérique peut détecter une erreur ; il ne démontre pas une identité pour tous les nombres.\end{attention}
\clearpage
\essSection{2}{Factoriser et choisir la forme utile}
\begin{definition}Factoriser transforme une somme en produit. On cherche un facteur commun : $6x^2-9x=3x(2x-3)$.\end{definition}
\begin{propriete}La différence de deux carrés se factorise :
\[a^2-b^2=(a-b)(a+b).\]
La double distributivité permet aussi d’établir
$(a+b)^2=a^2+2ab+b^2$ et $(a-b)^2=a^2-2ab+b^2$.\end{propriete}
\begin{exemple}$4x^2-25=(2x-5)(2x+5)$. Pour calculer $99^2$, écrire $(100-1)^2=10000-200+1=9801$.\end{exemple}
\begin{savoirfaire}[Prouver avec un programme de calcul]Nommer le nombre choisi $x$, traduire chaque étape, puis transformer l’expression. Pour « ajouter 4, multiplier par 3, retirer le triple du nombre initial », on obtient $3(x+4)-3x=12$ : le résultat est indépendant de $x$.\end{savoirfaire}
\begin{attention}$(a+b)^2$ n’est généralement pas $a^2+b^2$. Il faut les deux produits croisés, dont la somme vaut $2ab$.\end{attention}
\clearpage\rubrique{Découverte}
\begin{blocLivre}{ex-4-1}\begin{exercice}[1]{Remplacer une lettre}
Calculer $A=2x^2-3x+1$ pour $x=-2$ et pour $x=3$.
\end{exercice}\end{blocLivre}
\begin{blocLivre}{ex-4-2}\begin{exercice}[1]{Réduire une expression}
Réduire $7x-4+3x+9-2x$.
\end{exercice}\end{blocLivre}
\begin{blocLivre}{ex-4-3}\begin{exercice}[1]{Prendre l’opposé}
Supprimer les parenthèses et réduire $-(4x-7)+2(3-x)$.
\end{exercice}\end{blocLivre}
\begin{blocLivre}{ex-4-4}\begin{exercice}[1]{Double distributivité}
Développer et réduire $(3x+2)(x-5)$.
\end{exercice}\end{blocLivre}
\begin{blocLivre}{ex-4-5}\begin{exercice}[1]{Un facteur commun}
Factoriser $14x+21$ et $12x^2-8x$.
\end{exercice}\end{blocLivre}

\clearpage\rubrique{Entraînement}
\begin{blocLivre}{ex-4-6}\begin{exercice}[2]{Différences de carrés}
Factoriser $x^2-81$ et $25x^2-4$.
\end{exercice}\end{blocLivre}
\begin{blocLivre}{ex-4-7}\begin{exercice}[2]{Deux carrés de sommes}
Développer $(x+6)^2$ et $(2x-3)^2$.
\end{exercice}\end{blocLivre}
\begin{blocLivre}{ex-4-8}\begin{exercice}[2]{Calculer astucieusement}
Calculer $103\times97$ et $49^2$ en utilisant des identités.
\end{exercice}\end{blocLivre}
\begin{blocLivre}{ex-4-9}\begin{exercice}[2]{Factoriser en deux temps}
Factoriser $3x(x-2)+5(x-2)$, puis $9x^2-36$.
\end{exercice}\end{blocLivre}
\begin{blocLivre}{ex-4-10}\begin{exercice}[2]{Tester une conjecture}
Une élève affirme que $(x+2)^2-x^2=4$ pour tout $x$. Tester $x=1$, puis donner l’expression exacte.
\end{exercice}\end{blocLivre}

\clearpage\rubrique{Problèmes}
\begin{blocLivre}{ex-4-11}\begin{exercice}[3]{Deux programmes équivalents}
Programme A : ajouter 5 au nombre choisi, multiplier par 4, retirer 8. Programme B : multiplier le nombre choisi par 4, ajouter 12. Montrer qu’ils sont équivalents.
\end{exercice}\end{blocLivre}
\begin{blocLivre}{ex-4-12}\begin{exercice}[3]{Une bordure}
Un carré de côté $x$ cm est entouré d’une bordure de largeur 2 cm à l’extérieur. Exprimer l’aire de la bordure. Calculer pour $x=7$.
\label{fig-ex-4-12}\figureGeo{bordure}
\end{exercice}\end{blocLivre}
\begin{blocLivre}{ex-4-13}\begin{exercice}[3]{Des entiers consécutifs}
Montrer que le carré d’un entier impair est impair. Le carré d’un entier pair est-il toujours multiple de 4 ?
\end{exercice}\end{blocLivre}
\begin{blocLivre}{ex-4-14}\begin{exercice}[3]{Une égalité impossible ?}
Deux carrés ont pour côtés positifs $x+3$ et $x$. Leur différence d’aires peut-elle valoir 9 cm$^2$ ? Peut-elle valoir 39 cm$^2$ ?
\end{exercice}\end{blocLivre}
\begin{blocLivre}{ex-4-15}\begin{exercice}[3]{Une différence toujours positive}
Comparer, pour tout réel $x$, $A=(x+1)^2$ et $B=x(x+2)$. Démontrer la réponse.
\end{exercice}\end{blocLivre}

\clearpage
